% !TeX root = ../../main.tex
\chapter{稀疏矩阵计算的近似计算特征与性能瓶颈分析}\label{chap:analysis}
\section{问题提出}
稀疏性压缩了运算规模，却没有消除数据传输、索引解析及并行执行的不规则开销。当系统仍对所有数据使用相同数值精度，并按预先固定的范围执行计算时，应用输出所允许的简化空间往往没有得到利用。本章在具体系统设计之前，分析不同计算需求与实际执行成本之间的关系，为后续三章确定各自的优化对象。

本章结合三项研究的实验结果展开分析：利用 AccelES 的低比特候选覆盖实验考察阶段间精度需求，利用 ParetoES 的扫描比例与召回率曲线考察区域贡献差异，利用 EvoCG 的算子耗时与收敛轨迹考察状态相关的数值敏感性。分析重点是各类计算需求如何影响工作量和执行成本。由于实验平台、负载与质量指标不同，三项结果分别支撑对应的设计动机，性能比较限定在各自实验条件内。

\section{稀疏矩阵计算的基础特征}
\subsection{稀疏矩阵--向量乘法计算模型}
设 $A\in\mathbb{R}^{m\times n}$，$x\in\mathbb{R}^{n}$，矩阵第 $i$ 行非零元素的列索引集合为 $\mathcal{N}_i$。SpMV 可写为
\begin{equation}
 y_i=\sum_{j\in\mathcal{N}_i} A_{ij}x_j,\qquad 0\leq i<m.
 \label{eq:spmv}
\end{equation}
记非零元素总数为 $z=\nnz(A)$，通常将 $2z$ 作为一次 SpMV 的浮点操作量估计。该计数不包含索引运算、地址生成及排序等辅助工作，故只能用于描述算术工作量，不能单独预测执行时间。

在 Top-K SpMV 中，系统进一步返回
\begin{equation}
 \mathcal{T}_K=\TopK(Ax,K).
\end{equation}
这里 $\TopK$ 表示按得分选取索引集合；若存在并列值，应采用一致的索引优先规则。只有数据库向量和查询向量均经过适当的范数归一化时，内积排序才等价于余弦相似度排序。后续关于方向敏感聚类的讨论以归一化向量为基础，不能无条件推广到任意最大内积检索任务。

CG 求解 $Ax=b$ 时，要求 $A$ 为对称正定矩阵。其典型递推为
\begin{align}
 q_k&=Ap_k,&\alpha_k&=\frac{r_k^{\mathsf T}r_k}{p_k^{\mathsf T}q_k},\\
 x_{k+1}&=x_k+\alpha_kp_k,&r_{k+1}&=r_k-\alpha_kq_k,\\
 \beta_k&=\frac{r_{k+1}^{\mathsf T}r_{k+1}}{r_k^{\mathsf T}r_k},&p_{k+1}&=r_{k+1}+\beta_kp_k.
 \label{eq:cg}
\end{align}
因此，Top-K 中不被输出的行可能具有较低结果价值，而 CG 中某一局部乘积会通过递推影响后续状态，两类计算对近似误差的处理方式必须区分。

\subsection{稀疏数据的不规则访存特征}
以 CSR 为例，矩阵由行指针、列索引和非零数值构成。设数值与列索引分别占 $b_v$ 和 $b_c$ 字节，则仅矩阵主体读取量约为 $z(b_v+b_c)$，还需加入行指针、输入向量和输出向量的访问。若忽略缓存复用及额外开销，矩阵主体对应的算术强度近似为
\begin{equation}
 I_A\simeq\frac{2}{b_v+b_c}.
 \label{eq:intensity}
\end{equation}
例如，FP64 数值与 32 位列索引对应 $I_A\simeq 1/6$ FLOP/Byte。该估计描述仅计矩阵主体访问时的带宽压力。输入向量缓存不命中、额外索引或写回会进一步改变实际流量。

式\eqref{eq:intensity}还说明，降低数值位宽的收益受元数据比例限制。若数值缩小为原来的四分之一而索引保持不变，整个记录的字节数并不会同比减少。因此，低比特计算需要与索引压缩和总线打包联合优化。Ultra-CSR 的设计即针对数值压缩后行边界开销所占比例上升的问题。

不同访问阶段也要求不同布局。完整扫描适合连续传输，并能通过较长突发访问摊销启动开销；候选重算只访问少量离散行，要求快速定位行边界；簇选择则在簇间表现为不连续访问、在簇内表现为连续流式访问。数据结构若仅对一种模式优化，可能将算法减少的数据量重新转化为额外寻址和小粒度事务开销。

\subsection{稀疏计算的精度与并行执行特征}
降低精度可以同时影响存储占用、传输带宽及运算单元数量，但不同平台能够利用的位宽粒度不同。FPGA 可围绕 INT6 等自定义格式构建打包与乘法电路；GPU 则通常需要利用已有数据类型及线程束执行模型。因而，同样的量化策略在不同平台上的收益可能不同，不能只按位宽比例推算端到端加速。

稀疏结构与精度分布共同决定负载。按非零元素切分任务能缓解行长度不均，但需处理归约边界；按簇切分可保留局部性，但热点簇会造成执行拖尾；按 Tile 赋予不同精度后，各 Tile 的访问量与执行成本又可能不同。调度设计应使所选任务粒度与数据表示、归约机制及硬件资源相匹配。

\section{稀疏矩阵计算的近似计算机会分析}
\subsection{不同计算阶段的精度需求差异}
在百万行级检索任务中，最终返回的 $K$ 通常远小于数据库规模。若高精度计算的主要作用是确定少量候选的精确得分，则对显著落后于排序边界的大量行使用高精度，可能产生不必要的数据访问。低比特粗筛的目标应是保留真正有竞争力的候选，而不是逐个复现所有全精度得分。

\thesisfig[0.94]{acceles/Approximate_with_Top100_or_512.pdf}{低位宽与候选扩展对真实 Top-K 覆盖的影响}{fig:analysis-candidates}{%
图\ref{fig:analysis-candidates}展示 AccelES 的相关实验。直接以低比特 Top-100 代替真实 Top-100 可能产生明显偏差；将粗筛范围扩大到 Top-512 后，在展示的数据和查询条件下能够显著改善真实 Top-100 的覆盖。原论文比较 4、5、6 位表示，发现更低位宽可减少传输量，但也增加候选遗漏，说明位宽和候选预算需要联合选择\cite{jqzhai2025AccelES}。
}

为准确描述该现象，定义候选保持率

\begin{equation}
 R_{\mathrm{cand}}@K=\frac{|\mathcal{C}\cap\mathcal{T}_K|}{K},\qquad |\mathcal{C}|=L\geq K.
 \label{eq:candidate-recall}
\end{equation}
$\mathcal{C}$ 是近似预筛集合，$\mathcal{T}_K$ 是精确参考集合。当候选使用原始数值重算，并采用与参考一致的比较规则时，输出质量受候选覆盖程度限制。扩大 $L$ 往往降低遗漏风险，却增加重算和随机访问成本。因此，两阶段方案体现的是质量约束下的工作量分配，而非低位宽必然保持排序的理论保证。

\subsection{不同数据区域的计算贡献差异}
分层精度降低了全量扫描的单位成本，但仍需要读取整个低位宽数据库。如果查询相关候选集中于少量数据区域，则可在扫描前识别这些区域，进一步减少实际访问。关键在于建立低成本、具有足够辨别能力的区域代表，使区域选择与最终相似度目标一致。

\thesisfig[0.80]{paretoes/recall-scannedVectors_with_resparse.pdf}{聚类策略下候选扫描比例与检索召回率的关系}{fig:analysis-regions}{%
ParetoES 的实验将向量聚类后，考察扫描向量比例与 Recall@100 的关系，如图\ref{fig:analysis-regions}所示。部分数据上的曲线显示，召回率不会与扫描比例按固定比例增长，高相关区域能够贡献较多正确候选；不同聚类方法的曲线位置也存在明显差异。这一现象表明，方向相关性能够为结构感知选择提供依据，实际可缩减的访问范围取决于数据分布与召回目标\cite{paretoes2026}。
}

对于最终返回 $K$ 个候选的结果集合 $\widehat{\mathcal{T}}_K$，检索召回率为

\begin{equation}
 \mathrm{Recall}@K=\frac{|\widehat{\mathcal{T}}_K\cap\mathcal{T}_K|}{K}.
 \label{eq:recall}
\end{equation}
它与式\eqref{eq:candidate-recall}的分母相同，但被评估集合的大小和处理阶段不同。簇探测遗漏、元素剪枝、量化和局部候选截断均可能影响最终召回，因此端到端评估应覆盖完整处理链。

区域价值不均也意味着访问负载不均。频繁选中的簇会使对应计算核成为查询完成时间的决定因素。候选总量减少并不保证各核工作量同比减少；选择性计算必须同步解决热点分布和最慢计算核执行拖尾问题。

\subsection{不同运行状态的数值精度需求差异}
CG 的输出目标是求得满足残差要求的近似解，其误差含义与检索召回不同。以低精度执行第 $k$ 次 SpMV，可表示为
\begin{equation}
 \widetilde q_k=(A+E_k)p_k,
 \label{eq:perturbed}
\end{equation}
其中 $E_k$ 是对本次运算误差的等效描述。实际误差还与数值分布、缩放及累加方式有关。扰动经式\eqref{eq:cg}改变残差和搜索方向，既可能增加迭代次数，也可能影响最终可达精度。不精确 Krylov 方法提供了研究这一问题的理论框架\cite{simoncini2003,bouras2005}，但实际精度调度仍须与具体实现和求解目标结合。

图\ref{fig:analysis-cg-cost}给出 EvoCG 中 47 个矩阵按应用类别汇总的算子分析。SpMV 平均占求解时间的 67.1\%，表明其具有较大的优化价值；同时，其他算子、启动和归约仍然构成不可忽略的剩余开销。图\ref{fig:analysis-cg-traj}进一步显示，混合精度相对 FP64 在不同矩阵上的收敛轨迹并不相同：部分问题几乎不增加迭代，另一些问题会出现明显延迟\cite{evocg2027}。
\thesisfig[0.75]{evocg/spmv_breakdown.pdf}{CG 求解中 SpMV 的耗时占比}{fig:analysis-cg-cost}{%
SpMV 的平均耗时占比为67.1\%，说明数值流压缩具有较大的优化空间；其余时间由向量运算、归约及运行时开销构成，决定单一算子优化能够达到的端到端收益。
}

\thesisfig[1.00]{evocg/convergence.pdf}{混合精度与 FP64 求解的递推残差轨迹及精度变化}{fig:analysis-cg-traj}{%
上述差异说明，应结合运行状态决定精度，而不应预设所有矩阵均适用单调升精度或固定阈值策略。需要同时观察递推残差 $r_k$ 与用高精度重算的真实残差 $b-Ax_k$，避免将递推量达到阈值直接等同于真实解误差满足相同界限。本研究以递推残差、真实残差及收敛情况共同验证数值质量，结论适用于所评估的矩阵与求解条件。
}


\section{稀疏矩阵计算的主要性能瓶颈}
\subsection{高精度计算与数据传输开销}
两阶段计算的在线时间可以分解为
\begin{equation}
 T_{\mathrm{hier}}=T_{\mathrm{screen}}+T_{\mathrm{select}}+T_{\mathrm{refine}}+T_{\mathrm{transfer}},
 \label{eq:hiercost}
\end{equation}
其中各项分别对应低精度扫描、候选维护、高精度重算及接口传输。该式用于分解各阶段成本；存在流水重叠时，端到端时间由实际关键路径决定。低位宽减少第一项，但候选扩大可能提高后两项。数据格式必须同时降低连续扫描与离散候选读取的单位开销。

以 512 位传输为例，包含两个 32 位坐标和一个 FP32 值的 COO 记录占 96 位，一次传输最多容纳 5 个完整记录。若使用 10 位列索引、6 位数值和 1 位行边界标记，理论上可打包 30 个记录。该比较依赖列范围及非空行等假设，也没有计入所有控制数据；其意义在于说明数据格式压缩可以与低比特数值协同，提高单位传输的有效非零元素数量。

\subsection{全量数据访问与冗余计算开销}
选择性计算的成本包含簇代表比较、簇定位、簇内稀疏计算与最终结果合并。探测预算很小时，簇筛选和随机访问的固定开销可能占较大比例；探测预算接近簇总数时，又会逼近全量计算，并保留额外簇定位开销。因而，其优势取决于减少的扫描成本是否超过新增索引与调度成本。

对 $P$ 个并行核，若第 $p$ 个核处理的工作量为 $W_p$，在处理能力近似一致时，执行时间受 $\max_p W_p$ 而非平均工作量主导。本文使用变异系数
\begin{equation}
 \mathrm{CV}=\frac{\sqrt{P^{-1}\sum_{p=1}^{P}(W_p-\overline W)^2}}{\overline W}
\end{equation}
描述负载分散程度，并联合考察最大最小负载比及查询延迟。总体 CV 降低不必然保证每个单独查询都均衡，故还需在未见查询和分布偏移条件下评估。

\subsection{动态精度控制与异构执行开销}
对迭代求解，整体时间可以写为
\begin{equation}
 T_{\mathrm{solve}}=\sum_{k=0}^{N_{\mathrm{iter}}-1}
 (T_{\mathrm{SpMV},k}+T_{\mathrm{other},k}+T_{\mathrm{ctrl},k}).
 \label{eq:solvecost}
\end{equation}
精度策略同时影响单次 SpMV 时间和迭代次数。若忽略其他变化，单轮成本降为原来的 $a$ 倍、迭代次数增为 $b$ 倍，则总成本近似变为 $ab$ 倍，只有 $ab<1$ 才能产生加速。这是辅助解释的简化关系，实际分析应使用包含控制与运行时开销的实测时间。

EvoCG 在同一执行基础上平均增加 11.4\% 迭代而仍获得 1.28 倍几何平均加速，说明单轮节省可以超过迭代代价；但部分矩阵出现减速，说明这一条件并非普遍满足\cite{evocg2027}。策略推理、精度分支、数据类型转换以及多份数值存储也可能抵消收益。因此，需要固定结构索引、减少控制路径分歧，并降低策略进入迭代关键路径的频率和成本。

\section{面向近似计算的软硬件协同加速思路}
\subsection{低比特分层计算与高精度重算}
将候选识别与最终得分计算拆分，使第一阶段优先追求覆盖能力和扫描效率，第二阶段优先保证保留候选的得分精度。相应的数据格式与计算资源应按阶段专门设计：连续低比特数据流由多个并行核心处理，小规模高精度候选由可随机按行读取的重算路径处理。其适用条件包括候选集远小于数据库、候选保持率满足任务要求以及数据库可重复查询。

\subsection{结构感知选择性计算}
在低位宽执行基础上，以数据聚类建立候选区域，使查询只触发相关簇的读取和计算。系统应提供可配置探测预算、低成本簇选择及簇内连续访问，并根据代表查询的负载统计重排簇到计算核的映射。该方案适合存在可利用结构且允许召回率权衡的检索任务；对高召回要求或结构不可分的数据，其范围缩减空间会减小。

\subsection{状态感知动态混合精度计算}
当局部误差影响后续迭代时，需将精度作为与求解状态相关的动作。学习策略基于局部搜索方向作出精度选择，统一数据布局与 GPU 内核负责将异构精度转化为实际加速，运行时通过设备端执行和适度复用降低控制成本。训练可摊销性、矩阵分布变化和收敛验证是该方案的必要部署条件。

\begin{table}[htbp]\centering
\caption{三类近似计算问题的共同路径与不同边界}\label{tab:unified}
\begin{tabularx}{\textwidth}{p{2.25cm}XXX}\toprule
比较维度&低比特分层计算&结构感知选择性计算&状态感知混合精度\\\midrule
主要决策&粗筛位宽与候选数&访问哪些候选簇&局部区域使用何种精度\\
质量指标&候选保持率及重算结果&最终 Recall@K&残差、迭代次数与收敛性\\
主要执行挑战&流式与随机访问共存&簇访问热点与核间执行拖尾&精度异构与迭代控制开销\\
协同实现&双数据格式与异构核心&ACPE、微排序及热点均衡&Tile-CSC、统一内核及运行时\\
适用边界&预筛遗漏不能靠重算恢复&扫描比例与召回并无固定关系&策略迁移及收敛不具普遍保证\\\bottomrule
\end{tabularx}
\end{table}

% !TeX root = ../../main.tex
\section{计算需求差异的形式化分析}
\subsection{输出要求与近似预算的层次}
对计算任务定义输出要求，有助于确定哪些近似可以接受。检索的输入为固定数据库和查询，契约包括得分函数、目标 $K$、并列值处理以及召回阈值。CG 的契约则包括矩阵类别、右端项、初始解、停止准则和失败处理。两类契约中的阈值不能互换，例如检索 Recall 为0.99 并不能对应数值求解相对残差为0.01。

近似预算可以作用于多个层次。数值预算决定每个被处理元素使用多少位；范围预算决定处理多少候选或簇；控制预算决定多久重新选择精度。预算并非越小越好，因为减少一项开销可能增加另一项。例如更低位宽需要更大的候选集合，较少策略调用可能延迟必要的精度提升。分析应先明确预算影响的对象，再测量其对完整任务的影响。

\begin{table}[htbp]\centering\caption{不同预算的作用对象与补偿路径}
\begin{tabularx}{\textwidth}{p{2.2cm}XXX}\toprule
预算&直接节省&可能引入的代价&可用补偿方式\\\midrule
数值位宽&数值字段与算术资源&得分偏差、局部扰动&候选扩展、高精度重算\\
候选规模&精排读取与计算&相关条目被截断&扩大局部候选、重新筛选\\
簇探测数&数据库扫描量&相关区域被遗漏&提高探测预算、改善聚类\\
非零保留比例&非零记录数量&向量方向及得分变化&校准保留比例、保留原始副本\\
动作复用间隔&控制调用次数&响应状态变化滞后&缩短间隔、重新画像\\\bottomrule
\end{tabularx}\end{table}

\subsection{排序扰动的充分条件与候选扩展}
设近似分数为 $\hat y_i=y_i+e_i$，且 $|e_i|\leq\varepsilon_i$。真实 Top-K 中的行 $i$ 与集合外的行 $j$ 保持正确相对次序的充分条件是
\begin{equation}
 y_i-y_j>\varepsilon_i+\varepsilon_j.
\end{equation}
原因是近似差值的下界为 $\hat y_i-\hat y_j\geq y_i-y_j-\varepsilon_i-\varepsilon_j$。若右侧为正，近似次序不变。该条件关注排序间隔而不是分数本身的相对误差，因此两个具有相同平均误差的数据集可能拥有不同的候选召回。

将预筛规模从 $K$ 扩大到 $L$，需要关心真实 Top-K 是否被推到近似第 $L$ 名之后。对任意真实 Top-K 行，若最多只有 $L-1$ 个候选可能在扰动后超过它，则该行仍会保留。实际系统通常没有全部行的可用误差界，故通过验证查询估计覆盖。这一推导揭示了候选扩展的作用机制；实际配置的候选覆盖仍通过独立查询校准与验证。

当采用分布式局部候选时，风险又增加一层。即使全局近似 Top-$L$覆盖目标，局部分区容量不足仍可能删除重要行。质量审计应区分数值扰动和容量截断，分别使用“精确分数加局部选择”和“近似分数加全局选择”的对照，避免两种损失互相遮蔽。

\subsection{稀疏性与传输效率的分离}
逻辑稀疏性用 $z/(mn)$ 描述，而物理效率还取决于填充与事务粒度。设有效非零记录字节数为 $B_{\mathrm{valid}}$，实际读取为 $B_{\mathrm{actual}}$，则有效载荷率为
\begin{equation}
 \eta_{\mathrm{payload}}=B_{\mathrm{valid}}/B_{\mathrm{actual}}.
\end{equation}
剪枝使 $B_{\mathrm{valid}}$下降，但若每行仍需读取完整固定长度块，$B_{\mathrm{actual}}$可能不同比下降。类似地，过小的簇虽减少候选，却可能增加地址表和事务启动占比。

设每条总线记录包含数值 $b_v$、索引 $b_i$ 和边界信息 $b_f$ 位，总线宽度为 $W$，单包容量是
\begin{equation}
 P=\left\lfloor W/(b_v+b_i+b_f)\right\rfloor.
\end{equation}
该表达式是理想打包上限。实际实现还需考虑控制标记、跨包延续和对齐。比较不同格式时，应同时列出理论容量、实际包格式和测得带宽，避免把位级上限作为可持续性能。

\subsection{并行任务的不均衡与关键路径}
平均工作量只描述总体资源需求，完成时间还取决于任务最大值。若每核速率为 $v$，忽略其他开销，理想均衡时间为 $\sum_pW_p/(Pv)$，实际受最重核限制为 $\max_pW_p/v$。据此可定义不均衡因子
\begin{equation}
 \lambda=\frac{P\max_pW_p}{\sum_pW_p}\geq1.
\end{equation}
当工作量缩减比例为 $r$，不均衡因子却从 $\lambda_0$上升到 $\lambda_1$，简化模型下时间比例约为 $r\lambda_1/\lambda_0$。这说明选择性计算只有配合分配机制才可能稳定受益。

CV 与上述因子反映不同方面。CV 度量整体离散程度，最大值因子更接近同步查询的关键路径。两个负载集合可能 CV 相近而最大值不同。实验同时记录二者以及真实查询时间，才能区分平均分散改善与尾部拖延改善。

\subsection{预处理摊销与使用场景}
若某方案一次性准备时间为 $T_p$，单次执行时间为 $t$，连续执行 $N$ 次的总成本为 $T_p+Nt$。相对准备时间 $T_b$、执行时间 $t_b$ 的参考方案，只有
\begin{equation}
 N(t_b-t)>T_p-T_b
\end{equation}
才具有总时间优势。若 $t_b\leq t$，增加重复次数也不会弥补在线劣势。此关系适用于索引构建、数据转换和策略训练，但要求准备结果在这些执行中持续有效。

当数据库更新或矩阵变化使准备结果失效，重复次数应按两个更新事件之间的有效执行数计算，而不是系统生命周期内总请求数。尤其是 EvoCG 的训练策略可能与矩阵及求解设置相关，不能只因应用名称相同就认为训练成本永久摊销。

\subsection{综合评估的可比性条件}
不同实验应统一质量目标，却不必强行统一所有算法参数。相同探测数可能产生不同召回，相同位宽可能产生不同误差；因此，质量匹配优先于参数数值相同。为了隔离硬件差异，再补充固定算法和固定候选集合的对照。

\begin{table}[htbp]\centering\caption{跨方法比较中应固定与应报告的条件}
\begin{tabularx}{\textwidth}{p{2.5cm}XX}\toprule
比较目的&应固定条件&仍需报告的变化\\\midrule
近似算法收益&执行平台、输入、输出要求&位宽、候选预算与工作量\\
硬件架构收益&算法、候选集合、质量要求&数据布局、并行映射与接口开销\\
精度控制收益&稀疏内核与运行时基础&精度轨迹、迭代数、策略开销\\
负载均衡收益&同一查询序列与总任务&映射、峰值核负载、数据迁移\\
部署成本收益&有效使用周期与更新频率&预处理、功耗范围、复用次数\\\bottomrule
\end{tabularx}\end{table}

上述框架将质量约束、工作量变化与执行成本联系起来。本文实验覆盖其中与三项研究直接相关的条件组合，跨平台、跨任务的结果则分别在各自评估范围内解释。统一负载下的综合验证仍有待进一步开展。

% !TeX root = ../../main.tex
\section{从计算需求到执行成本的分层分析}
\subsection{算子、阶段与应用三个计时边界}
同一项优化在不同计时边界下可能产生不同结论。算子边界只覆盖 SpMV 或 Top-K 选择；阶段边界包含输入准备、算子调用和结果整理；应用边界还可能包含数据编码、索引构建与下游消费。本文主要讨论查询与求解的在线完整流程，并单独列出离线成本。明确边界的意义在于，只有实际落在关键路径中的成本变化才会影响用户可见的完成时间。

以两阶段检索为例，若仅测量 INT6 扫描内核，所观测到的是候选分数生成速率；若测量一次查询，还需要包含局部候选维护、合并与 FP32 重算。低比特扫描越快，其他阶段在总时间中的相对比例越高。因此，不能用低比特算子的加速比直接代表查询加速比。若不同阶段对连续查询形成流水，系统吞吐又由最慢阶段决定，单条查询延迟则仍包含其经过各阶段的时间。

对 CG，单轮 SpMV 用时下降并不充分。设参考实现执行 $N_0$ 轮，每轮平均时间 $t_0$；近似实现执行 $N_1$ 轮，每轮核心计算平均时间 $t_1$，每轮控制平均时间为 $t_c$，其他一次性在线成本为 $T_s$。则简化的求解加速比为
\begin{equation}
 S=\frac{N_0t_0}{N_1(t_1+t_c)+T_s}.
\end{equation}
该式没有假定每轮时间完全一致，而是以完整轨迹的平均量表示总和。它揭示两种约束：一是额外迭代抵消单轮节省，二是短任务中一次性成本的比例较大。真正的性能目标因而应是满足相同质量条件的总求解时间。

\begin{table}[htbp]\centering\caption{不同计时边界能够支持的结论}
\begin{tabularx}{\textwidth}{p{2.4cm}XX}\toprule
边界&可回答的问题&不能直接推出的结论\\\midrule
SpMV 内核&单位稀疏计算的执行效率&完整查询或求解的加速比\\
候选生成阶段&粗筛成本及候选覆盖情况&精排后的全部任务质量\\
在线查询&既定索引上的查询代价&包含索引更新的长期服务成本\\
完整迭代求解&达到停止条件的时间&包含策略训练的首次求解收益\\
有效使用周期&准备、更新与执行的总成本&不同更新模式下的同等收益\\\bottomrule
\end{tabularx}\end{table}

\subsection{数值位宽与实际传输量的关系}
考虑一个用于说明字段成本的理想记录：数值占 $b_v$ 位，列索引占10 位，暂不计行边界及对齐。将数值从32 位降到6 位，记录从42 位变为16 位，其字节数比例为 $16/42$，与 $6/32$ 不同。若再把512 位传输单元视为不可分割的包，容量分别为12 项与32 项，比例又受到取整影响。这些数值是字段预算示例，不替代第3 章的实际 Ultra-CSR 协议。

实际传输还存在固定开销。设有 $z$ 个保留非零元素，元数据总字节数为 $M$，每个元素平均物理字节数为 $b$，则
\begin{equation}
 B=M+zb.
\end{equation}
量化改变 $b$，剪枝改变 $z$，结构压缩可能改变 $M$。当 $M$ 占比很高时，继续减少数值位宽的边际收益会降低；当矩阵由大量极短行组成时，行指针、边界标记和包尾空闲尤其需要关注。仅报告数值数组大小会低估这些影响。

还应区别总传输量和传输速率。Random-CSR 只访问少量候选，但各行起点不连续，单次读取长度可能较短；Ultra-CSR 访问更多记录，却能形成长连续流。因此，两种路径的有效带宽并不必然相同。合理的两阶段设计允许高带宽流承担大量低成本筛选，让低带宽随机流只服务少量必要重算，而不是假定所有数据读取具有相同成本。

\subsection{工作量减少与任务碎片化的关系}
选择性计算把原本连续扫描拆为若干被选区域。假设一次查询访问 $s$ 个簇，第 $j$ 个簇有效数据为 $B_j$，每次启动的固定成本为 $\tau_j$，其有效传输速率为 $v_j$，串行路径上的读取成本可写为
\begin{equation}
 T_{\mathrm{read}}\approx\sum_{j=1}^{s}\left(\tau_j+\frac{B_j}{v_j}\right).
\end{equation}
该式不代表 HBM 控制器的精确排队模型，但说明相同总字节数可以对应不同延迟。把一个区域拆得更细，能够提高选择的针对性，也可能增加启动次数和地址处理。若这些事务可以重叠，固定成本的一部分被隐藏；是否隐藏则取决于并发请求能力、缓存和依赖关系。

因此，簇数增加具有双重作用：一方面改善候选区域的分辨率，另一方面增加代表评分、目录与短事务开销。聚类目标若只优化簇内相似度，就未必得到最快的物理执行。第4 章通过簇内连续布局和多核探测将这两种要求联系起来，但它仍需要根据具体数据分布选择工作点。

GPU 混合精度也存在类似现象。过细的动作粒度可以减少高精度元素，却可能产生更多小任务和更多调度信息。统一内核的作用是将异构数值选择限制在共享的结构遍历中，减少把动作差异放大为内核碎片的机会。由此可见，近似决策的粒度必须与硬件任务粒度共同设计。

\subsection{质量损失来源的逐层定位}
检索系统中的低召回可能来自量化、元素剪枝、簇选择、局部 Top-K 容量以及并列分数规则。若只报告最终 Recall，无法判断应该调整哪一层预算。可采用逐层恢复的分析方式：保留其余执行不变，仅把某层替换为精确或全量版本，观察结果变化。这里的“精确”必须指定参考数值类型与得分函数，而不是把任何浮点计算称为数学精确。

\begin{table}[htbp]\centering\caption{检索质量损失的定位方法与局限}
\begin{tabularx}{\textwidth}{p{2.2cm}XX}\toprule
可能来源&可采用的对照&对照仍不能消除的因素\\\midrule
数值量化&保留候选范围，改用 FP32 评分&范围遗漏和局部容量截断\\
非零剪枝&保留格式，恢复完整非零支撑&量化和选择机制的影响\\
簇级选择&保持评分方式，探测全部簇&数值扰动与局部截断\\
局部候选上限&汇集全体近似分数做全局选择&近似分数本身的排序扰动\\
结果并列规则&统一稳定行号作为次级键&其他误差来源\\\bottomrule
\end{tabularx}\end{table}

这些对照之间可能存在交互，因而各自改善的 Recall 不能简单相加。某一候选同时被簇选择和量化排除，单独修复其中一层时仍无法恢复它。完整分析需要观察候选集合的交、并与遗漏原因，而不只是报告几个百分比差值。

CG 的定位方式不同。需要在相同内核和相同运行时下比较精度策略，并保存残差轨迹、精度轨迹与每轮时间。对某轮发生的残差变化，可回放相同搜索方向分别执行 FP64 与混合精度 SpMV，分析局部扰动；但离线回放不包含扰动引起的后续状态改变，不能替代完整求解。两种分析分别回答单步误差和闭环后果，作用互补。

\subsection{从共性认识到章节技术选择}
上述分析表明，三项研究可以共享问题分解方式，但不能共享未经证明的性能或误差结论。第3 章选择较粗的阶段分工，因为候选得分可以在筛选后重新计算；第4 章选择簇级范围控制，因为数据区域具有可利用的相关性；第5 章选择反馈精度控制，因为当前误差会改变后续求解状态。

三者也具有不同的撤销能力。候选进入重算集合后，可以恢复其高精度分数，但被提前删除的候选无法由精排自动恢复。CG 中已经发生的低精度更新通常不能靠下一轮简单升精度完全撤销，可能需要残差校正或重启。由此决定了质量检查位置和保守预算不同，不能将检索中的高精度重算直接类比为迭代中的无代价纠错。

从执行角度看，FPGA 能够围绕固定数据协议配置算术位宽与排序逻辑，适合将阶段结构固化为流水；GPU 则依靠共享内核和批量执行吸收动态动作，适合在通用并行资源上控制运行时开销。本文依据任务语义选择硬件执行方式，使决策粒度、存储布局和并行组织相互匹配。

\section{本章小结}
本章围绕阶段精度、候选区域和迭代状态三类需求差异，分析了稀疏矩阵计算中的近似机会及其执行代价。分析表明，降低位宽、缩减候选范围与动态调整精度均可减少部分运算和访存，但实际收益同时受到元数据、任务不均衡和控制开销的制约。因此，近似决策需要与数据表示、并行粒度和执行机制协同设计。候选保持、检索召回与数值收敛则分别限定各类方法的质量边界，为后续三章的方法设计与实验评估提供依据。

在这一框架下，检索实验以一致的得分定义、目标集合和质量要求比较查询吞吐；迭代求解实验以一致的停止条件和真实残差检查比较求解时间。通过同时考察质量、决策开销及完整执行路径，后续研究进一步区分局部计算节省与端到端性能改善之间的关系。
